%! AP Statistics — Unit 8, Lesson 8.2 — Guided Notes
\documentclass[10pt]{article}
\usepackage{apstats-article}
\usepackage{apstats-boxes}

\begin{document}

\pageheader{Unit 8, Lesson 8.2}{Guided Notes}
\namedateperiod

\begin{objectivebox}
By the end of this lesson, I will be able to:
\begin{itemize}
    \item[$\square$] \writeline
    \item[$\square$] \writeline
    \item[$\square$] \writeline
    \item[$\square$] \writeline
    \item[$\square$] \writeline
\end{itemize}
\end{objectivebox}

\begin{vocabbox}
\termblanklong{Chi-square distribution}
\termblanklong{Goodness of fit (GOF) test}
\termblanklong{Expected count}
\termblanklong{Null proportion ($p_0$)}
\termblanklong{Degrees of freedom ($df$)}
\termblanklong{Large Counts condition}
\end{vocabbox}

\begin{hookbox}
A die is rolled 120 times and the results are recorded. The claimed distribution is
that each face appears with probability $1/6$.

Why can't we use six separate one-proportion $z$-tests to check this? \writelines{2}
\end{hookbox}

\begin{notesbox}{The Chi-Square Distribution (VAR-8.A)}
\textbf{Properties:}
\begin{itemize}
    \item Values are always \blank{4cm} (never negative).
    \item Shape is skewed \blank{3cm}.
    \item As degrees of freedom \blank{4cm}, the skew becomes less pronounced.
    \item $df = (\text{number of categories}) - $ \blank{1cm} for a GOF test.
\end{itemize}

\vspace{0.1in}
\textbf{When to use a GOF test (VAR-8.C):} Use a chi-square GOF test when you have
\blank{3cm} categorical variable and want to test whether the population follows
a \blank{5cm} distribution.
\end{notesbox}

\begin{notesbox}{Stating Hypotheses (VAR-8.B)}
H$_0$ states the \blank{5cm} for \blank{4cm} category.

H$_a$: At least \blank{3cm} of the proportions is not as specified in H$_0$.

\vspace{0.1in}
\textbf{Die-Roll Example:}

H$_0$: $p_1 = $ \blank{1.5cm}, $p_2 = $ \blank{1.5cm}, $p_3 = $ \blank{1.5cm},
$p_4 = $ \blank{1.5cm}, $p_5 = $ \blank{1.5cm}, $p_6 = $ \blank{1.5cm}

H$_a$: \writeline

\textit{Note:} All null proportions must sum to \blank{1cm}.
\end{notesbox}

\begin{notesbox}{Calculating Expected Counts (VAR-8.D)}
Formula: \quad $E = $ \blank{4cm} $\times$ \blank{4cm}

\vspace{0.1in}
\textbf{Die-Roll Example} ($n = 120$): \quad $E = 120 \times \dfrac{1}{6} = $ \blank{3cm} per face.

\vspace{0.1in}
\begin{center}
\renewcommand{\arraystretch}{1.4}
\begin{tabular}{lcccccc}
    \textbf{Face} & 1 & 2 & 3 & 4 & 5 & 6 \\
    \hline
    \textbf{Null proportion} & $1/6$ & $1/6$ & $1/6$ & $1/6$ & $1/6$ & $1/6$ \\
    \textbf{Expected count} & \blank{1.5cm} & \blank{1.5cm} & \blank{1.5cm}
                            & \blank{1.5cm} & \blank{1.5cm} & \blank{1.5cm} \\
\end{tabular}
\end{center}
\end{notesbox}

\begin{notesbox}{Verifying Conditions (VAR-8.E)}
\begin{enumerate}
    \item \textbf{Random:} The data come from a \blank{5cm} or randomized experiment.
          If sampling without replacement: $n \le $ \blank{2cm}$\% \cdot N$.
    \item \textbf{Large Counts:} All \blank{4cm} counts are $\ge $ \blank{1cm}.
\end{enumerate}

\vspace{0.1in}
\textbf{Die-Roll Check:}
\begin{itemize}
    \item Random: \blank{8cm}
    \item Large Counts: Smallest expected count is \blank{2cm}. Condition \blank{3cm}.
\end{itemize}

\vspace{0.05in}
\textit{Important:} The Large Counts condition uses \blank{4cm} counts, not
\blank{4cm} counts.
\end{notesbox}

\begin{practicebox}
\textbf{Practice:} A geneticist expects offspring to appear in a 9:3:3:1 phenotype
ratio. She observes 160 offspring. How many does she expect in each phenotype category?
Show the formula and fill in the table.

\vspace{0.1in}
Total ratio parts: $9 + 3 + 3 + 1 = $ \blank{2cm}

\begin{center}
\renewcommand{\arraystretch}{1.4}
\begin{tabular}{lcccc}
    \textbf{Phenotype} & A & B & C & D \\
    \hline
    \textbf{Null proportion} & $9/16$ & $3/16$ & $3/16$ & $1/16$ \\
    \textbf{Expected count} ($n=160$) & \blank{2cm} & \blank{2cm} & \blank{2cm} & \blank{2cm} \\
\end{tabular}
\end{center}

Is the Large Counts condition met? Explain. \writelines{1}
\end{practicebox}

\end{document}
